% Section 10.2, from lectures/L10/appendix/b_exercises.md: the description of app::EchoNode, and
% the exercise that builds it.

\section{What You Are Building: \code{app::EchoNode}}
\label{c10:sec:echo_node}\label{c10:app:b}

\subsection{The \code{app::EchoNode} application}
\label{c10:sec:echo_app}
\code{EchoNode} is the application that finishes the book: it receives a byte over the UART and
sends it straight back. On the bench, what you type in the terminal comes back to you, having
crossed every layer the book built. The application itself is tiny, because all the hard work is
already done and tested underneath it: the driver, the transport, the peripheral. \code{EchoNode}
just spends that stack.

The important idea is that \code{EchoNode} is \textbf{software you write and host-test before you
touch the bench}. It depends only on \code{driver::uart::Interface}, not on \code{Uart},
\code{AvrSpi}, or any hardware, so you build it, test it against the UART stub, and only \emph{then}
flash it and bring it up, as the last rung of the ladder (Exercise~\ref{c10:ex:c3}). By the time it
reaches the bench, its logic is already proven; the only thing the bench adds is the wiring. This is
the same discipline the whole book runs on --- host-test behind the interface seam first, hardware
last --- applied one level up, to the application.

This section \textbf{describes} \code{EchoNode}; you write it yourself. The UART stub it is tested
against was written in \chapref{6}, as soon as the interface it implements existed.

\subsection{The application seam: \code{app::Interface}}
\label{c10:sec:app_interface}
Applications sit behind their own interface, \code{app::Interface}, exactly as drivers sit behind
\code{driver::uart::Interface}. It declares a single operation, \code{run(const volatile bool& stop)}, which
runs the application until \code{stop} becomes \code{true}. \code{main()} holds an
\code{app::Interface&} and calls \code{run()}, so it neither knows nor cares which application it
runs; swapping \code{EchoNode} for another is a one-line change in \code{main()}.

\code{stop} is a \code{const volatile bool&}, a flag the application only \emph{reads}, hence the
\code{const}, owned by whoever started it. Setting it \code{true} asks the application to return from
\code{run()}, which is a clean shutdown and, in a test, the way the loop is ended. It is
\code{volatile} because it is written outside the loop that reads it: by the UART stub in a test,
and on a target by whatever else may set it, an interrupt handler for instance. \code{run()} never
writes the flag, and where nothing the loop calls could write it either, as with an interrupt
handler, the compiler may read it once, keep it in a register and never see it change;
\code{volatile} makes every pass read it from memory. It is a \code{volatile bool} rather than a
\code{std::atomic}, because avr-libc is freestanding and ships no \code{<atomic>}, and on the
single-core ATmega a byte read is already indivisible: atomicity is not the problem here, visibility
is.

\subsection{\code{app::EchoNode}}
\label{c10:sec:echo_class}
\code{EchoNode} implements \code{app::Interface}. Write it as \file{include/app/echo_node.hpp} (the
declarations) and \file{source/app/echo_node.cpp} (the definitions), in the \code{app} namespace, in
the same AVR-portable style as the driver (\code{<stdint.h>} and bare types, nested namespace
blocks, no \code{[[nodiscard]]}). It is \code{final}, non-copyable, and non-movable.

\textbf{Member variable.} There is one: \code{driver::uart::Interface& myUart}, the UART driver to
echo over, injected through the constructor and held by reference. That reference is what lets the
\emph{same} class run over the concrete \code{Uart} on hardware and over the UART stub in a test,
with no change. Because it is a reference it is set once at construction and never rebound, which is
also why copy and move are deleted: a reference member cannot be reassigned, the same reason
\code{Uart} deletes them.

\textbf{Methods.} The constructor, \code{EchoNode(driver::uart::Interface& uart)}, is
\code{explicit} and \code{noexcept}. It takes the driver to echo over and stores it in \code{myUart},
and it does no I/O, since construction only wires the dependency and therefore cannot fail.

\code{run(const volatile bool& stop)} is the application loop, \code{noexcept}, overriding
\code{app::Interface::run()}. It repeats until \code{stop} is \code{true}, and on each pass it checks
\code{stop}, asks \code{myUart} for a byte with the \textbf{non-blocking} \code{read()}, and, if one
arrived, echoes it straight back with \code{writeBlocking()}.

The receive is a poll rather than \code{readBlocking()} on purpose. A blocking read would wait
forever for a byte and never look at \code{stop} again, so the loop could not be stopped, and a test
of it would hang. Polling \code{read()} returns to the top of the loop every pass to re-check
\code{stop}, which is exactly what lets a caller, or the test, end it. The echo write may block,
using \code{writeBlocking()}, because the byte is already in hand and simply has to go out.

\subsection{Host-testing it, before any hardware}
\label{c10:sec:echo_test}
Because \code{EchoNode} depends only on \code{driver::uart::Interface}, you test it against the
\textbf{UART stub} you wrote back in \chapref{6} (Exercise~\ref{c6:ex:4.1}), the moment that
interface existed: no \code{Uart}, no transport, no FPGA. This is where it finally earns its keep.
The trick that makes \code{run()} return in a test is the \code{stop} flag: the stub sets
\code{stop} to \code{true} the moment its scripted RX runs out, so \code{run()} echoes every queued
byte and then ends on the next check. The test queues the input, constructs the \code{EchoNode} over
the stub, calls \code{run(stop)}, and asserts that the stub's recorded TX equals what it queued, in
order. The suite lives in \file{test/app/echo_node_test.cpp} and, like the other host tests, may use
full modern C++; run it with \code{make test}.

That termination trick is also a check on your implementation: because it depends on \code{run()}
polling \code{read()} and re-checking \code{stop}, a \code{run()} that blocked in
\code{readBlocking()} would hang the test --- the test telling you the loop is not actually
stoppable.

\subsection{On the bench}
\label{c10:sec:echo_bench}
Nothing about the class changes for hardware. The application \code{main()} on the Nano constructs
the real stack under the same interface --- an \code{AvrSpi}, a \code{Uart} over it,
\code{configure()}d for 115200~8N1, then an \code{EchoNode} over the \code{Uart} --- and calls
\code{run()} with a \code{stop} flag it leaves \code{false}, so the node echoes forever. That is the
last rung of the bring-up ladder (Exercise~\ref{c10:ex:c3}): the exact class you host-tested, now
echoing real characters typed in the terminal, over real SPI, through your own VHDL peripheral.

\begin{aside}
\textbf{How to check your work.} The provided host suite \file{fw/test/app/echo_node_test.cpp}, run
by \code{make test}, checks \code{EchoNode} over the UART stub; it switches on once
\file{app/echo_node.hpp} and \file{driver/uart/stub.hpp} both exist. The answers to the parts that
are not code are in \secref{sol:c10}. Try each part before you read its answer.
\end{aside}

\exercise{\code{app::EchoNode}}{C++}
\label{c10:ex:b1}

\xpart{a} Write \file{include/app/interface.hpp} first: the \code{app::Interface} seam described
above, with a virtual \code{noexcept} destructor that is \code{= default}, and the single pure
virtual \code{run(const volatile bool& stop) noexcept}. Then write \file{include/app/echo_node.hpp} and
\file{source/app/echo_node.cpp} from the description above --- the \code{myUart} member, the
constructor, \code{run(const volatile bool& stop)}, and the deleted copy and move --- so that \code{EchoNode}
implements \code{app::Interface}. Run \code{make test} and confirm that the echo test passes with no
hardware.

\xpart{b} The test drives \code{run()} with the stub, which stops once its scripted RX is exhausted,
queuing the bytes \code{0x00}, \code{0x41} and \code{0xFF}. Add a case that queues \textbf{nothing}
and confirm that \code{run()} returns immediately, having sent nothing. Which implementation mistake
is this the only case that catches?

\xpart{c} The method \code{run(const volatile bool& stop)} is the unit under test, ended by the flag. Explain
why \code{run()} must poll the non-blocking \code{read()} rather than call \code{readBlocking()}, and
connect it to why \code{stop} is passed by reference (and read every pass) rather than returned or checked once, and why
it is \code{volatile}.

\xpart{d} Change \code{EchoNode} to echo each byte back \textbf{uppercased} (leave non-letters
alone), and update the test. Keep the transformation inside \code{run()}; the point is that the
application layer is exactly where such policy belongs, above a driver that only moves bytes.

\subsection{What's ahead}
With \code{EchoNode} written and host-tested, \secref{c10:sec:exercises} is the bench: flash it and
climb the bring-up ladder to watch it echo on real hardware, then trace one byte through every layer
of the stack the book built.
